% TOP_PROBLEM: {"id": 8700028, "title": "Conjecture 3.10 — Gel’fond-Schneider", "category_id": 3, "category": {"id": 3, "name": "graph_theory", "display_name": "Graph Theory", "description": "Problems involving graphs, networks, and their properties.", "slug": "graph-theory", "order_index": 3, "created_at": "2026-07-31T15:26:25.671Z"}, "status": "open", "problem_number": "AMR-086-0028", "set_id": 13, "statusQualification": "Restores the algebraic-input, exact-degree, and common nonzero-logarithm hypotheses from the primary source."}
% FIRST_POSTED: 2026-10-01
% ENTRY_KIND: statement_only
% UnsolvedMath ID 8700028; source catalogue code AMR-086-0028
% !TeX program = lualatex
% Definition and sources only; this file is not an LLM solution attempt.
\documentclass[11pt,a4paper]{article}
\usepackage[margin=25mm]{geometry}
\usepackage{fontspec}
\setmainfont{lmroman10-regular.otf}[BoldFont=lmroman10-bold.otf,
 ItalicFont=lmroman10-italic.otf,BoldItalicFont=lmroman10-bolditalic.otf]
\usepackage{amsmath,amssymb,mathtools}
\usepackage{unicode-math}
\setmathfont{latinmodern-math.otf}
\AtBeginDocument{\let\setminus\smallsetminus}
\usepackage{xurl,hyperref}
\hypersetup{colorlinks=true,urlcolor=blue}
\setcounter{secnumdepth}{0}
\setlength{\parindent}{0pt}
\setlength{\parskip}{5pt}
\setlength{\emergencystretch}{3em}
\begin{document}
\section*{Conjecture 3.10 — Gel’fond-Schneider}\label{8700028}
{\small UnsolvedMath ID 8700028 · AMR-086-0028 · Graph Theory · AMR Open Problem Lists}
\subsection{Definitions and mathematical statement}
Let $\alpha\in\overline{\mathbb Q}^{\times}$, and choose a nonzero complex logarithm $\lambda$ with $e^\lambda=\alpha$. Let $\beta$ be algebraic irrational of degree $d=[\mathbb Q(\beta):\mathbb Q]\ge2$. With this one logarithm fixed, define $\alpha^z=e^{z\lambda}$ for every complex $z$.

The combined Gelfond--Schneider conjecture asks whether the following $d$ complex numbers are algebraically independent over $\mathbb Q$:
\[
\lambda,\quad e^{\beta\lambda},\quad e^{\beta^2\lambda},
\quad\ldots,\quad e^{\beta^{d-1}\lambda}.
\]
Explicitly, evaluating any nonzero polynomial in $d$ variables with rational coefficients at this ordered list should give a nonzero complex number. Equivalently, the field they generate over $\mathbb Q$ should have transcendence degree $d$.

The quantification is over every permitted algebraic pair and every nonzero logarithm of $\alpha$. No principal-branch restriction is imposed. The possibility $\alpha=1$ is included when the chosen logarithm is nonzero. These input conditions belong to the original source's setup and are needed to make the catalogue's abbreviated statement complete.

The power in the $j$th exponential is $\beta^j\lambda$, so all entries must use the same $\lambda$. The conclusion simultaneously excludes relations among the exponential entries and relations also involving the logarithm. Proving algebraic independence after omitting $\lambda$, or proving independence of each pair containing $\lambda$, would not by itself establish the full assertion for general $d$.
\subsection{Short English statement}
Are a chosen nonzero algebraic logarithm and all the indicated powers formed from it jointly algebraically independent?
\subsection{Sources}
\begin{itemize}
\item[S1] ULAM.AI and the UnsolvedMath Contributors, supplied problems.json, record 8700028 (AMR-086-0028), AMR Open Problem Lists. Catalogue identity and source transcription; CC BY 4.0.
\url{https://huggingface.co/datasets/ulamai/UnsolvedMath}.
\item[S2] Waldschmidt - Open Diophantine Problems (2004), Conjecture 3.10. Original source indexed by AMR-086-0028.
\url{https://arxiv.org/abs/math/0312440}.
\item[S3] M. Waldschmidt, Open Diophantine Problems, Moscow Mathematical Journal 4 (2004), 245–305. Author PDF supplies the surrounding definitions and hypotheses.
\url{https://webusers.imj-prg.fr/~michel.waldschmidt/articles/pdf/odp.pdf}.
\end{itemize}
\end{document}
